%% ================================================================
%%  NEW SECTION DRAFT — FAILURE MODE ANALYSIS AND
%%  HALLUCINATION TAXONOMY
%%
%%  Intended placement: between §5 Experimental Evaluation and §6
%%  Design Insights.  LaTeX auto-numbering will renumber subsequent
%%  sections.
%%
%%  Status: DRAFT — awaiting v6 benchmark results for final numbers
%%  in §X.1 and empirical validation of class-level guardrail impact
%%  in §X.4.  All v5 statistics below are from
%%  benchmark/full_run_v5_improvements.log and benchmark/summary.txt
%%  of the v5 run committed in bd3a0cb.
%%
%%  2026-04-11 addendum: §X.6 Operational Failure Modes (O-1..O-3)
%%  added at supervisor's request after the 2026-04-11 22:28 incident
%%  where a v6 rerun aborted at 90 s on "Credit balance is too low"
%%  hidden inside CLI stdout.  The three pipeline countermeasures
%%  cited in §X.6 are implemented in commit 811028c.
%% ================================================================

\section{Failure Mode Analysis and Hallucination Taxonomy}
\label{sec:taxonomy}

Section~\ref{sec:results} reports Pinpoint's aggregate accuracy against
the baselines.  Before we turn to design insights
(§\ref{sec:insights}), we dissect \emph{why} Pinpoint fails when it
fails.  Prior work on VLM-based GUI grounding treats grounding errors
monolithically---Nguyen's iterative narrowing paper, for example,
describes its own failure cases as ``context loss''~\cite{nguyen2024iterativenarrowing}
without further decomposition.  We find that Pinpoint's errors are
neither monolithic nor random: they decompose into a small number of
structurally distinct hallucination classes, each rooted in a specific
prompt-design or perceptual property.  Naming these classes is the
prerequisite for targeted interventions.

\subsection{A Bimodal Error Distribution}
\label{sec:taxonomy-bimodal}

Re-examining the 50-case pre-guardrail benchmark (Section~\ref{sec:experiments})
case by case rather than in aggregate, we find that the Pinpoint error
distribution is not unimodal.  Two disjoint regimes dominate:

\paragraph{Mode~A ($34/50 = 68\%$, median $29$\,px).}
When the Phase~1 grid search lands within the Phase~2 search radius
$s_0 = 100$\,px of the true target, Pinpoint's Monte Carlo dot
convergence succeeds: the median error in this subset is
$29$\,px, and \emph{two cases reach $1.0$\,px}---confirming that
pixel-level localization is achievable with a single VLM family
when the algorithm is on-target.

\paragraph{Mode~B ($16/50 = 32\%$, median $155$\,px).}
When Phase~1 lands outside $s_0$ of the true target, Pinpoint does not
recover.  Phase~2's per-round ``HIT'' on a hex-stencil dot simply
selects the dot that looks closest to the VLM's internal expectation
of the target within the already-misplaced search disc, returning a
garbage answer with full confidence.  Critically, the catastrophic
failures
($>$200\,px) in this regime are \emph{tightly correlated} between
Grid and Pinpoint---for example, on \texttt{ui\_dense\_08} a Grid
estimate $388$\,px from the target is followed by a Pinpoint
estimate $391$\,px from the same target.  Phase~2 does not correct
Phase~1 errors; it faithfully propagates them.

The overall Pinpoint median of $57$\,px across all $50$ cases is
therefore a mixture artefact.  It describes neither Mode~A's
near-target convergence ($29$\,px) nor Mode~B's cascade-propagated
breakdown ($155$\,px); it is the centroid of two populations that
require different interventions.  A Pinpoint that merely sharpens
Mode~A will not move the aggregate median by much as long as Mode~B
dominates the right tail.

\subsection{Anatomy of a Mode~B Failure: Cascading Hallucination}
\label{sec:taxonomy-cascade}

Tracing through a single Mode~B case makes the mechanism concrete.
The target is \texttt{ui\_button\_03}: ``the centre of the `Open'
button in the toolbar'', ground truth $(514, 39)$.  Phase~1's
axis-wise grid search lands at $(279, 21)$---it has correctly
localized the \emph{row} of toolbar buttons (within $18$\,px
vertically) but picked the wrong button within that row, a
$235$\,px horizontal error that places the search disc entirely
away from the target.

Phase~1.5 then asks the VLM to declare the anchor.  The prompt is
shown the Phase~1 context crop, which now contains the wrong button
but not the target.  It is instructed to ``describe the exact pixel
location you will target''.  The target is not in the crop, but the
prompt does not allow the VLM to say so: it is a declarative command,
not a binary check.  The VLM resolves the conflict in the only way it
can---by re-interpreting the target to fit what it sees.  It declares
the anchor as \emph{``the exact center pixel of the purple `Open'
button''}, describing a visually similar toolbar button that happens
to occupy the centre of the Phase~1 crop.

Every subsequent round of Phase~2 now uses this corrupted description
as its judgement criterion.  The majority vote
(Section~\ref{sec:phase2}) faithfully picks the dot closest to
the \emph{described} anchor, which is the wrong button's centre.
After eight rounds of convergence, Pinpoint returns $(261, 24)$---
$253$\,px from the true target---with unanimous agreement across
all voting rounds.

The failure is not random; it is \emph{cascading}: a single
declaration error---forced by a prompt that does not permit
non-existence---propagates through eight rounds of self-reinforcing
confirmation.  A majority vote is powerless because every voter is
ranking against the same contaminated anchor.

This walkthrough motivates our taxonomy.  The cascading failure is
not one bug but a chain: each step exposes a different hallucination
class, each with its own prompt-level or perceptual root cause.

\subsection{Hallucination Classes (H-$\ast$)}
\label{sec:taxonomy-classes}

We identify seven visual-grounding hallucination classes (H-1
through H-7) by analyzing the exact prompts, image rendering,
and voting procedure of our pre-guardrail Pinpoint implementation.  The
classes are summarized in Table~\ref{tab:taxonomy} and elaborated
below.  A given failure
typically involves multiple classes; Mode~B failures are dominated by
\textbf{H-3}.

\begin{table*}[t]
  \centering
  \small
  \caption{A taxonomy of VLM grounding hallucinations observed in
    the pre-guardrail Pinpoint.  Severity reflects the frequency with which the
    class appears in benchmark failures and its contribution to
    Mode~B.  ``Root'' marks the class identified as the dominant
    cause of cascading Mode~B failures.}
  \label{tab:taxonomy}
  \begin{tabular}{@{}llp{6.2cm}p{3.5cm}c@{}}
    \toprule
    \textbf{ID} & \textbf{Name} & \textbf{Failure mechanism} &
    \textbf{Prompt / render trigger} & \textbf{Severity} \\
    \midrule
    H-1 & Forced choice         & Prompt insists the VLM pick one dot,
                                   producing false-positive HITs when
                                   the target is absent                       & ``you must pick one''          & High \\
    H-2 & Pseudo majority       & Identical prompts at temperature~0
                                   yield non-independent votes, so a
                                   ``majority'' filter behaves like a
                                   single sample                              & 3 votes of same image, $T{=}0$  & Medium \\
    H-3 & Anchor contamination  & Declaration forced on an absent-target
                                   crop; VLM picks a similar element as
                                   anchor, corrupting all later rounds        & ``declare the target pixel''   & \textbf{Root} \\
    H-4 & Target occlusion      & Filled dots physically cover the very
                                   pixels they purport to mark, denying
                                   the VLM ground evidence                   & radius-2 filled circles         & Medium \\
    H-5 & Salience capture      & Bright dot colours and numeric labels
                                   dominate attention at the expense of
                                   the underlying target                     & 7 saturated hues                & Low--Med \\
    H-6 & Label displacement    & Numeric labels drawn at $+$3,$-$10
                                   offset from dot centre; VLM may
                                   reference the label, not the dot          & text offset relative to marker  & Low \\
    H-7 & Relative ``closest''  & ``closest'' is scale-free: any dot can
                                   be closest, even 50\,px away               & no distance threshold           & Medium \\
    \bottomrule
  \end{tabular}
\end{table*}

\paragraph{H-1 Forced choice.}
The Phase~2 prompt includes the phrase ``there is always a closest
dot---you must pick one''.  This wording was introduced to discourage
the VLM from refusing to answer, but it also eliminates the
possibility of the VLM reporting that no dot is on the target.  In
Mode~B cases, where the search disc does not contain the true target,
the VLM is structurally compelled to return a dot.  The returned
number is indistinguishable, from the caller's perspective, from a
genuine HIT.

\paragraph{H-2 Pseudo-majority vote.}
\label{par:h2-pseudo-majority}
The pre-guardrail implementation calls the VLM three times per Phase~2 round on
\emph{identical} images and prompts at temperature zero, and takes
the majority answer.  With a deterministic decoder this is not a vote
over independent samples: all three calls are expected to return the
same answer, and any systematic bias (e.g.\ colour preference, label
position bias) appears in all three votes in the same direction.  The
majority filter therefore removes sporadic decoding noise but cannot
correct structural bias.  This observation is the empirical caveat to
condition~(ii) of Theorem~\ref{thm:convergence}'s convergence
guarantee: the independent-trial assumption underlying
Table~\ref{tab:convergence-error} is violated under this decoder
configuration.

\paragraph{H-3 Anchor contamination (root cause of Mode~B).}
This is the class exhibited in §\ref{sec:taxonomy-cascade}.  The
anchor declaration prompt is declarative, not interrogative.  It
assumes the target is visible and asks the VLM to describe its exact
pixel.  When Phase~1 is off-target, the VLM has no way to push back:
it must produce a description, and it produces the best description
it can of whatever \emph{is} in the crop.  This description then
becomes the judgement criterion for eight rounds of Phase~2, each of
which faithfully ranks dots against it.  Faithfulness bias
\cite{kadavath2022language} guarantees that the VLM will find a
winner.  Because the anchor is the only object that persists across
rounds, corrupting it corrupts everything downstream.

\paragraph{H-4 Target occlusion.}
Dots are rendered as 4-pixel filled circles with a 1-pixel white
outline.  Their purpose is to mark candidate positions, but a
filled dot placed on the target necessarily \emph{covers} the
target's pixels in the image the VLM sees.  The VLM is then judging
``is this dot on the target?'' using only the pixels at the dot's
perimeter---a far weaker evidence source than the pixel actually
under the dot centre.  This is a second-order effect but compounds
H-5 and H-7.

\paragraph{H-5 Salience capture.}
At the resolution the VLM perceives, the dots are some of the most
visually salient features in the image: small, saturated, distinct
from almost any background.  This is deliberate for legibility, but
it creates a risk that VLM attention becomes dominated by the dots
themselves rather than their relationship to the target.  Anecdotally
we have observed the VLM's self-reported reasoning refer to dot
colour properties (``the red dot'') more often than target properties
(``the button label''), suggesting dot-centric rather than
target-centric attention.

\paragraph{H-6 Label displacement.}
Each dot is drawn with its numeric label at offset $(+3, -10)$ from
the dot centre, so the label does not overlap the dot.  The VLM is
instructed to ``identify which numbered dot is closest to the target''
and reply with ``D\textless number\textgreater''.  But the spatial
reference for ``D3'' is ambiguous: is it the location of the coloured
circle, or the location of the text ``3''?  A 10-pixel upward offset
is within one image patch at typical vision-encoder resolutions, so
the VLM's notion of ``dot~3's location'' may be biased several pixels
upward.  The effect is small but systematic.

\paragraph{H-7 Relative ``closest''.}
The phrase ``which dot is \emph{closest}'' defines closeness as a
ranking, not a threshold.  If every dot is 50\,px away from the
target, the VLM will still return ``the closest one''.  There is no
prompt-level mechanism for the VLM to say ``no dot is \emph{on} the
target''.  Combined with H-1, this means that Phase~2 rounds never
report ``MISS for the right reason''.

\subsection{Guardrails as Structured Responses}
\label{sec:taxonomy-guardrails}

With the failure classes named, the guardrail design problem becomes
concrete: each guardrail should be analyzable as a targeted response
to one or more classes, and the set of guardrails should cover the
set of classes.  Table~\ref{tab:guardrail-matrix} maps the guardrails
implemented or proposed in this work to the classes they address.

\begin{table}[t]
  \centering
  \small
  \caption{Guardrails (rows) by hallucination class (columns).
    $\bullet$ = primary target, $\circ$ = secondary effect.  P1 and
    P2 (G-1, G-2) are implemented in the guardrail-applied run reported in
    Section~\ref{sec:results}; P3 (G-3), textual labelling (G-5),
    perturbation voting (G-4), hollow rings (G-6), and directional
    bisection (G-7) are future work targeting visual grounding
    (H-1..H-7).}
  \label{tab:guardrail-matrix}
  \resizebox{\columnwidth}{!}{%
  \begin{tabular}{@{}lccccccc l@{}}
    \toprule
    \textbf{Guardrail} & \multicolumn{7}{c}{\textbf{Visual grounding (H-1..H-7)}} & \textbf{Status} \\
    \cmidrule(lr){2-8}
                       & H-1 & H-2 & H-3 & H-4 & H-5 & H-6 & H-7 & \\
    \midrule
    G-1 Visibility check (P2)              & $\bullet$ &           & $\bullet$ &           &           &           & $\circ$   & applied \\
    G-2 Zoom cap 1024\,px (P1)             &           &           &           & $\circ$   &           &           &           & applied \\
    G-3 Hybrid post-HIT verification (P3)  & $\circ$   &           &           &           &           &           & $\bullet$ & future \\
    G-4 Perturbation majority              &           & $\bullet$ &           &           & $\circ$   &           &           & future \\
    G-5 Letter labels (A, B, C$\dots$)     &           &           &           &           & $\bullet$ & $\circ$   &           & future \\
    G-6 Hollow ring dots                   &           &           &           & $\bullet$ &           &           &           & future \\
    G-7 Directional bisection              & $\bullet$ &           &           & $\bullet$ &           & $\bullet$ &           & future \\
    \bottomrule
  \end{tabular}%
  }
\end{table}

\paragraph{G-1: visibility check (this work, P2).}
We replace the declarative anchor prompt (``describe the exact target
pixel'') with an interrogative one (``is this target clearly visible?
VISIBLE:\textellipsis\ or ABSENT:\textellipsis'').  This gives the VLM
a structured escape hatch before it commits to a description, breaking
the H-3 faithfulness chain at its source.  On an \texttt{ABSENT}
response the anchor declaration is retried at a wider crop (one third
of the screen) and, if still absent, on the full image.  Persistent
absence is logged as a Mode~B signal for downstream confidence
scoring.  G-1 also closes H-1 because the VLM is no longer in a
``pick one'' regime; and it partially addresses H-7 because the
``visible'' condition is a de facto proximity threshold.

\paragraph{G-2: zoom cap (this work, P1).}
Independent guidance from both the Claude API optimization documentation and
empirical investigation indicates that Claude~3.5~Sonnet silently downscales
images whose longer side exceeds roughly $1092$\,pixels at 1:1 aspect
ratio.  Our pre-guardrail pipeline upscales context crops to $1200 \times 1200$
and refinement crops to $1600 \times 1600$---both in the downscaled
regime, losing approximately~17\% of the spatial resolution the
operator believes they are providing.  Capping the upscale at $1024
\times 1024$ eliminates this silent loss at no cost in latency or
queries.  G-2 is a secondary contribution to H-4: sharper pixel
evidence at the dot's edge.  We note that the $1024$-pixel cap
is VLM-resolution-dependent: it was tuned to Sonnet~4's
effective resolution.  For higher-resolution VLMs such as
Opus~4.7 ($\sim\!2576$~px), the cap should scale accordingly;
the principle generalizes, the constant does not
(Section~\ref{sec:resolution-method}).

\paragraph{Future guardrails (G-3 through G-7).}
The remaining guardrails are deferred to future work but merit brief
discussion because the taxonomy motivates them precisely.
\textbf{G-3} adds a zero-shot post-HIT verification call---``is the
target within 5\,px of dot D\textit{n}?''---that must use a completely
independent API call with no prior conversational history to avoid
confirmation bias.  \textbf{G-4} replaces the identical-input
majority vote with input perturbation: cyclic colour shift and
horizontal mirroring with inverse transform, so that a systematic
bias in any single rendering is broken in the others.  \textbf{G-5}
replaces numeric ``D3''-style labels with single-letter labels
``A, B, C, \textellipsis'', exploiting the VLM's OCR pathway which
prior work reports as $\approx$20\% more accurate at spatial
localization than colour-based labels~\cite{you2023ferret}.
\textbf{G-6} uses hollow ring dots with transparent centres so that
the VLM can see the target pixels under a dot.  \textbf{G-7}
replaces the closest-dot question with a directional bisection
(``is the target north or south of dot A?''), converting an $N$-way
classification into a sequence of binary relative-spatial
judgements.  VLMs show measurably higher log-probability separation
for binary spatial questions than for $N$-way categorical
ones~\cite{seeclick, showui}.

\subsection{Discussion}

Two observations from this taxonomy are worth highlighting.

\paragraph{Root cause is prompt shape, not model capability.}
H-3, the class we identify as the root cause of Mode~B, is not a
limitation of the underlying VLM's visual reasoning.  The same model
that produces a contaminated anchor in an absent-target crop is able,
given the same crop and an \emph{interrogative} prompt, to correctly
report that the target is not present---as our G-1 experiment
(Section~\ref{sec:results}) demonstrates.  The failure was built into
the prompt.  This is a hopeful finding: it suggests that the ceiling
on VLM grounding accuracy is not yet bounded by model capability but
by the design of the prompting pipeline that wraps it.

\paragraph{Aggregate metrics hide mode structure.}
A median of $57$\,px and a mean of $87$\,px, in isolation, would lead
a researcher to conclude that VLM grounding needs ``better models''
or ``more training data''.  The bimodal structure reveals a very
different diagnosis: in 68\,\% of cases Phase~1 lands within the
Phase~2 search radius and Pinpoint converges to a median of $29$\,px
(with two cases at $1.0$\,px), while in the remaining 32\,\% Phase~2
cascades from a misplaced search disc into a garbage answer whose
median distance is $155$\,px.  The failure in the latter 32\,\% is
not an intrinsic precision limit; it is an \emph{entrained} error
whose root cause is a single forced-commitment prompt.  We therefore
encourage future VLM grounding papers to report
\emph{mode-stratified} accuracy---accuracy conditioned on whether the
coarse search landed in the correct neighbourhood---alongside
aggregate metrics.  Publishing only aggregates can mask
qualitatively different failure populations and misdirect research
effort toward the wrong bottleneck.

\subsection{Operational Failure Modes}
\label{sec:operational-failures}

The seven visual-grounding classes H-1 through H-7 describe failures
\emph{internal to the model's reasoning}: the VLM is fed a well-formed
image and prompt, and the resulting answer is wrong.  A parallel
family of failures, which we denote with the O-$\ast$ family,
occurs \emph{outside} the model---in the surrounding pipeline that
shuttles images, prompts, and answers between components.  Operational
failures are not cognitive errors; they are infrastructural faults
that the model itself cannot observe, and that aggregate metrics
silently absorb unless the pipeline is instrumented to detect them.
We document them here because a grounding system that reports 58\,\%
Mode~B in the field is indistinguishable from one reporting 58\,\%
O-1 in the benchmark logs: both look like catastrophic failure, but
the remedies are disjoint.  This distinction was forced on us by a
real incident during this work; we report it as a worked example.

\paragraph{O-1 Cascading silent failure.}
The VLM API call returns a non-zero exit status, but the surrounding
code path does not distinguish ``error'' from ``empty answer''.  The
benchmark's \texttt{\_vlm\_call\_cli} helper captures \texttt{stdout}
with \texttt{subprocess.run(capture\_output=True)}, and the underlying
\texttt{claude\ -p} command writes its error message to
\texttt{stdout} rather than \texttt{stderr}.  As a result, the caller
logs only a generic \texttt{CLI\ error\ (rc=1):} line, swallows the
specific cause (in our incident, \emph{``Credit balance is too
low''}), and proceeds to the next case.  Every subsequent case
records a \texttt{FAIL} result that is indistinguishable from a
genuine grounding failure.  In a long run this can consume hours of
wall clock time and produce a results file that passes sanity checks
but contains no signal.

\paragraph{O-2 Phantom miss via rate limit.}
The API returns an HTTP success status but with a throttled or
truncated body.  The caller parses a blank response as an empty
answer, which the benchmark records as \texttt{FAIL} with distance
\texttt{null}.  Unlike O-1 the exit status is clean, so no warning is
emitted; the failure only becomes visible when a per-category failure
rate spikes above its historical baseline.  A fleet of such phantom
misses can halve headline accuracy without tripping any single
alarm.

\paragraph{O-3 Credit exhaustion mid-run.}
An otherwise healthy run is interrupted part-way because a metered
resource (API credit, disk quota, session lease) is depleted.  If the
pipeline writes results in a single batch at the end of the run, all
partial progress is lost on interrupt.  This occurred in our own pre-guardrail
benchmark: 115 minutes of data were lost when the run was terminated
near its cost ceiling, because \texttt{results.json} was written only
after the outer loop completed.

\paragraph{O-4 Audit-trail invisibility.}
A different gap appears at the boundary between a guardrail and the
operators who must trust its decisions.  When a mid-pipeline
guardrail silently rejects, quarantines, or transforms its
input---a dual-LLM quarantine boundary, an output sanitization
step, a policy filter---downstream consumers observe only the
post-filter result.  Operators investigating a false positive or a
false negative have no structured trail to reconstruct what the
guardrail saw, which mitigation fired, or why.  This is distinct
from O-2 phantom miss, where an external API silently returns a
throttled empty body: in O-4 the guardrail acted exactly as
designed, but its action was not externalised in a queryable form.
Aggregate accuracy figures absorb both classes into the same FAIL
bucket while their remedies are disjoint.  Standard structured
logging only partially addresses this, because a fully internalised
log couples the engine to a persistence layer the engine has no
business knowing about.  The mitigation we landed in this work---an
injected callable audit sink, with each guardrail emitting a
Pydantic-typed row (\texttt{LLMQAuditEntry} and
\texttt{SanitizationAuditEntry}) through it---preserves the
engine-server boundary while admitting opt-in observability,
mirroring the inversion-of-control pattern previously used for
approval-gate callbacks in the same engine.  The API was added at
commit \texttt{4464902} on 2026-05-07; production wiring of the
sink to a durable store remains future work.

\paragraph{Countermeasures.}
Three pipeline changes render O-1 through O-3 survivable, and a
state-callback audit sink renders O-4 reconstructible, all without
any change to the model or prompts.
\textbf{First}, a \emph{90-second smoke check}: within 90--120 seconds
of launch, the harness verifies that the first test case produced a
non-null prediction.  If not, the run is aborted and the operator is
notified.  In our incident this bounded the damage of O-1 to a single
case---instead of several hours.
\textbf{Second}, \emph{JSONL partial writes with fsync}: each case
writes one line to a \texttt{results.json.partial} file and flushes
to disk before the next VLM call.  A corrupted trailing line (from
interrupt mid-write) is tolerated on reload.  This renders O-3
fully recoverable.
\textbf{Third}, \emph{unbuffered stdout}: replacing the default
block-buffered stream with \texttt{line\_buffering=True} makes the
live log file reflect real-time progress, so an operator watching a
tail can catch O-1 and O-2 long before the smoke check would.
\textbf{Fourth}, a \emph{structured audit emission API via
state-callback sink}: each privileged guardrail emits Pydantic-typed
audit rows---one schema per guardrail class---through a callable
injected into workflow state at startup; the engine carries no
compile-time dependency on the persistence layer, and a
downstream server-side consumer can attach the sink to a durable
store.  This renders O-4
detectable by reconstruction without coupling the engine to log
infrastructure, and the JSON round-trippable schema admits standard
log aggregator interop.
These three pipeline changes were introduced in our codebase as a
single commit after the incident; we recommend them as defaults for
any benchmark harness that wraps a metered external model.  The
fourth, an admission rubric, was added subsequently at the workflow
policy layer as agent ecosystems expanded beyond a single model
boundary.  The fifth, an opaque-key-anchored polling primitive, was
implemented directly in the dialogue server used in this work as a
coordination-layer primitive that renders the cross-channel race
condition unrepresentable by construction.  The sixth, a
rephrase-confirm protocol, was adopted as a coordination norm
during a development session in which two parse-divergence
incidents occurred within hours, and is applied at the dispatch
layer rather than the codebase.  The seventh, a structured audit
emission API, was implemented in the engine layer at commit
\texttt{4464902} as part of a pattern-coverage integration phase,
with downstream persistence wiring staged for production.

Operational failures occupy a different corner of the failure space
than hallucinations, but they enter the same numerator in any
headline accuracy figure.  A grounding system evaluated only on
end-to-end accuracy cannot tell whether it lost 10\,\% to H-3 or to
O-1, and the two require orthogonal remedies.  We therefore suggest
that benchmark reports separate the two axes explicitly: one number
for cognitive failures (H-$\ast$) and one for operational failures
(O-$\ast$), with the latter expected to trend toward zero as
infrastructure matures.

%% ================================================================
%%  END NEW SECTION DRAFT
%% ================================================================
